
\subsection*{Order-Preserving and Order-Reversing Behavior}

\begin{remark*}[Order-reversing terminology]
The orderings chapter defines an \hyperref[def:order-reversing-map]{order-reversing map}.
Such maps are also called \emph{antitone}. This section uses that existing
definition and studies how order-preserving and order-reversing maps interact
with injectivity, restriction, composition, and inverse behavior.
\end{remark*}

\begin{dependencies}
\begin{itemize}
  \item \hyperref[def:order-preserving-map]{Order-Preserving Map}
  \item \hyperref[def:order-reversing-map]{Order-Reversing Map}
  \item \hyperref[def:function]{Function}
\end{itemize}
\end{dependencies}

\begin{definitionbox}{Definition (Strictly Monotone Map)}
\begin{definition}[Strictly Monotone Map]\label{def:strictly-monotone-map-functions-order}
Let $(M,\leq)$ and $(M',\leq')$ be ordered sets, and write $x<y$ for
$x\leq y$ and $x\neq y$.
\begin{enumerate}[label=(\roman*)]
\item $f:M\to M'$ is \emph{strictly order-preserving} if
$x<y$ implies $f(x)<' f(y)$.
\item $f:M\to M'$ is \emph{strictly order-reversing} if
$x<y$ implies $f(y)<' f(x)$.
\end{enumerate}
\end{definition}
\end{definitionbox}

\begin{dependencies}
\begin{itemize}
  \item \hyperref[def:order-preserving-map]{Order-Preserving Map}
  \item \hyperref[def:order-reversing-map]{Order-Reversing Map}
  \item \hyperref[def:strict-partial-order]{Strict Partial Order}
\end{itemize}
\end{dependencies}

\begin{remark*}[Comment]
In many analysis texts, ``monotone'' can mean increasing or decreasing, and
``strictly monotone'' means strictly increasing or strictly decreasing. In this
chapter, plain monotone means order-preserving unless the order-reversing case
is named explicitly.
\end{remark*}

\begin{propositionbox}{Proposition (Constant Maps Are Monotone)}
\begin{proposition}[Constant Maps Are Monotone]\label{prop:constant-maps-monotone}
Let $(M,\leq)$ and $(M',\leq')$ be ordered sets. Every constant function
$f:M\to M'$ is order-preserving.
\hyperref[prf:constant-maps-monotone]{Go to proof.}
\end{proposition}
\end{propositionbox}

\begin{dependencies}
\begin{itemize}
  \item \hyperref[def:constant]{Constant Function}
  \item \hyperref[def:order-preserving-map]{Order-Preserving Map}
  \item \hyperref[def:reflexive]{Reflexive Relation}
\end{itemize}
\end{dependencies}

\begin{remark*}[Comment]
The preceding proposition is the standard warning: monotone does not imply
injective. If $M$ has two distinct elements, a constant map is monotone but
not injective.
\end{remark*}

\begin{theorembox}{Theorem (Strict Monotonicity Implies Injectivity on Linear Orders)}
\begin{theorem}[Strict Monotonicity Implies Injectivity on Linear Orders]\label{thm:strict-monotone-injective-linear}
Let $(L,\leq)$ be a linear order and let $(P,\leq')$ be a partial order.
If $f:L\to P$ is strictly order-preserving or strictly order-reversing, then
$f$ is injective.
\hyperref[prf:strict-monotone-injective-linear]{Go to proof.}
\end{theorem}
\end{theorembox}

\begin{dependencies}
\begin{itemize}
  \item \hyperref[def:strictly-monotone-map-functions-order]{Strictly Monotone Map}
  \item \hyperref[def:totally-ordered-set]{Totally Ordered Set}
  \item \hyperref[def:injective]{Injective Function}
\end{itemize}
\end{dependencies}

\begin{remark*}[Comment]
If a strictly monotone map on a linear order is not injective, then two
comparable distinct inputs must have been sent to the same output. That
contradicts strictness, because strictness demands a strict inequality between
their outputs.
\end{remark*}

\subsection*{Composition Rules}

\begin{theorembox}{Theorem (Composition Rules for Monotonicity)}
\begin{theorem}[Composition Rules for Monotonicity]\label{thm:monotone-composition-rules}
Let $f:A\to B$ and $g:B\to C$ be functions between ordered sets.
\begin{enumerate}[label=(\roman*)]
\item If $f$ and $g$ are order-preserving, then $g\circ f$ is order-preserving.
\item If $f$ and $g$ are order-reversing, then $g\circ f$ is order-preserving.
\item If exactly one of $f$ and $g$ is order-reversing and the other is
order-preserving, then $g\circ f$ is order-reversing.
\end{enumerate}
\hyperref[prf:monotone-composition-rules]{Go to proof.}
\end{theorem}
\end{theorembox}

\begin{dependencies}
\begin{itemize}
  \item \hyperref[def:composition]{Composition}
  \item \hyperref[def:order-preserving-map]{Order-Preserving Map}
  \item \hyperref[def:order-reversing-map]{Order-Reversing Map}
\end{itemize}
\end{dependencies}

\begin{remark*}[Comment]
The first item expands to
\[
a_1\leq a_2
\Longrightarrow
f(a_1)\leq' f(a_2)
\Longrightarrow
g(f(a_1))\leq'' g(f(a_2)).
\]
The other cases differ only by the direction of the second inequality.
\end{remark*}
